\section*{Chapter \ref{ch:iv1:linear-algebra} --- Linear Algebra}
\markboth{\protect\omTocMark{}{Answers: Linear Algebra}}{\protect\omTocMark{}{Answers: Linear Algebra}}

\begin{solution}{iq:iv1:linear-algebra:1}
The matrix $uv^\top$ has rank 1, and its eigenvalues are $v^\top u = 6$, 0 and 0. The matrix of all ones has eigenvalues $n$ and $0$ ($n - 1$ times). With zeros on the diagonal, they are $n - 1$ and $-1$ ($n - 1$ times).

For any vector $x$, $uv^\top x = u(v^\top x)$\index{rank-one matrices}, a multiple of $u$. So the image is one line, and $u$ is an eigenvector: $uv^\top u = (v^\top u)u$. Everything orthogonal to $v$ goes to 0. The trace is $v^\top u$, as a check.

The matrix of all ones is $\mathbf{1}\mathbf{1}^\top$, with $\mathbf{1}^\top\mathbf{1} = n$. Taking away the identity shifts each eigenvalue by $-1$.

\iqwrong{Computing a characteristic polynomial. Spot the structure first: a shift by the identity moves all the eigenvalues, and it leaves the eigenvectors alone.}
\end{solution}

\begin{solution}{iq:iv1:linear-algebra:2}
The eigenvalues are $2.8$ once, with eigenvector $\mathbf{1}$, and $0.8$ nine times, on everything orthogonal to $\mathbf{1}$. The lowest $\rho$ is $-1/9$, and the variance of the average falls to $\rho$ for large numbers.

Write the matrix as $(1 - \rho)I + \rho\mathbf{1}\mathbf{1}^\top$\index{equicorrelation}. The matrix of all ones has eigenvalues $n$ and 0, so the eigenvalues here are $1 - \rho$ and $1 - \rho + n\rho = 1 + (n - 1)\rho$.

A correlation matrix needs all eigenvalues at or above zero, so $\rho \ge -1/(n - 1)$. At that value, the sum of the assets is constant.

The average has variance $\frac{1}{n^2}\mathbf{1}^\top C\mathbf{1} = \frac{1 + (n - 1)\rho}{n} = \rho + \frac{1 - \rho}{n}$. Here it is $0.28$. This is the limit of diversification: the shared part never goes away.

\iqwrong{Thinking that any $\rho$ in $[-1, 1]$ is possible. With many assets, they cannot all be strongly negatively correlated with each other.}
\end{solution}

\begin{solution}{iq:iv1:linear-algebra:3}
No. Its determinant is $-0.62$, and the portfolio $(1, -\sqrt{2}, 1)$ has a variance of $4 - 3.6\sqrt{2} \approx -1.09$, which is impossible.

A valid correlation matrix must be positive semi-definite\index{positive semi-definite matrices}, because $w^\top Cw$ is the variance of a portfolio. All 2 by 2 blocks look fine, so looking at pairs does not catch this. By symmetry, try $w = (1, -a, 1)$. Its variance is $1 + a^2 + 1 - 2a(0.9) - 2a(0.9) = 2 + a^2 - 3.6a$. That is lowest at $a = 1.8$, and already negative at $a = \sqrt{2}$.

Intuitively, if $X$ is close to $Y$ and $Y$ is close to $Z$, then $X$ cannot be uncorrelated with $Z$.

\iqwrong{Checking that each number lies in $[-1, 1]$. Those are necessary conditions on pairs, but the requirement is on the whole matrix.}
\end{solution}

\begin{solution}{iq:iv1:linear-algebra:4}
The eigenvalues can only be 0 or 1. The rank equals the trace, and the trace of $X(X^\top X)^{-1}X^\top$ is $p$.

If $Pv = \lambda v$, then $P^2v = \lambda^2v$, and so $\lambda^2 = \lambda$\index{projection matrices}. Since $P(P - I) = 0$, the polynomial $x(x - 1)$ has distinct roots, so $P$ is diagonalisable. The trace is the count of 1s, and the rank is also the count of 1s.

For the hat matrix, $\operatorname{tr}\bigl(X(X^\top X)^{-1}X^\top\bigr) = \operatorname{tr}\bigl((X^\top X)^{-1}X^\top X\bigr) = \operatorname{tr}(I_p) = p$. The cyclic rule does the work.

\iqwrong{Assuming that a projection must be symmetric. The matrix with rows $(1, 1)$ and $(0, 0)$ is idempotent and not symmetric, and its rank and trace still agree.}
\end{solution}

\begin{solution}{iq:iv1:linear-algebra:5}
The eigenvalues are 3 and 1, with eigenvectors $(1, 1)$ and $(1, -1)$. So $A^{10}$ has diagonal $(3^{10} + 1)/2 = 29\,525$ and off-diagonal $(3^{10} - 1)/2 = 29\,524$, and $A^{-1}$ has diagonal $2/3$ and off-diagonal $-1/3$.

The matrix is $I + \mathbf{1}\mathbf{1}^\top$, so its eigenvalues follow at once\index{spectral decomposition}. It is 3 times one projection plus 1 times another, and the product of these two projections is zero. So any function of $A$ applies to the 3 and to the 1 only:
\[
f(A) = \frac{f(3)}{2}\begin{pmatrix}1&1\\1&1\end{pmatrix} + \frac{f(1)}{2}\begin{pmatrix}1&-1\\-1&1\end{pmatrix}.
\]
With $f(x) = 1/x$, the diagonal is $(\tfrac13 + 1)/2 = \tfrac23$, and the off-diagonal is $(\tfrac13 - 1)/2 = -\tfrac13$. Direct multiplication gives the same.

\iqwrong{Raising the entries to the power 10. Powers of a matrix act on its eigenvalues, not its entries.}
\end{solution}

\begin{solution}{iq:iv1:linear-algebra:6}
The correlation between $X$ and $Z$ can be anywhere from $0.62$ to 1.

From the determinant: a 3 by 3 correlation matrix is valid when\index{correlation bounds}
\[
1 + 2abc - a^2 - b^2 - c^2 \ge 0,
\]
which, for a given $a$ and $b$, gives $c \in \bigl[ab - \sqrt{(1 - a^2)(1 - b^2)},\ ab + \sqrt{(1 - a^2)(1 - b^2)}\bigr]$. Here that is $0.81 \pm 0.19$.

From angles: a correlation is the cosine of the angle between two random variables seen as vectors. $\arccos 0.9 \approx 25.8$ degrees. Angles between $X$ and $Z$ can be from 0 to $2 \cdot 25.8$ degrees, and the cosine at the top end is $\cos(2\arccos 0.9) = 2 \cdot 0.81 - 1 = 0.62$.

Correlation is not transitive in general. With correlations of $0.1$ and $0.1$, the third can be anywhere from $-0.98$ to 1, which is no real constraint. A strong constraint appears only when the correlations are high.

\iqwrong{Multiplying $0.9 \cdot 0.9 = 0.81$, and giving that as the answer. It is the centre of the interval, and it would be the answer if $X$ and $Z$ were linked only through $Y$. They need not be.}
\end{solution}

\begin{solution}{iq:iv1:linear-algebra:7}
The eigenvalues of $P$ are the $n$-th roots of unity, the eigenvectors are the Fourier vectors $(1, \omega, \omega^2, \dots)$ and the condition number is 1. The circulant has eigenvalues $3$, $1$, $2 + i$ and $2 - i$.

If $x_k = \omega^k$, with $\omega^n = 1$, then $(Px)_k = \omega^{k + 1} = \omega x_k$\index{circulant matrices}. Since $P$ is orthogonal, $P^\top P = I$, all its singular values are 1, which gives the condition number.

A circulant with first row $(c_0, c_1, \dots)$ is the polynomial $c_0I + c_1P + \dots$ in $P$. So it has the same eigenvectors, and its eigenvalues are $c_0 + c_1\omega + \dots$. Here, $2 + \omega$ with $\omega \in \{1, i, -1, -i\}$ gives $3$, $2 + i$, $1$ and $2 - i$.

\iqwrong{Saying that a real matrix has real eigenvalues. Only symmetric ones must. A shift is a rotation in disguise, and a rotation has complex eigenvalues.}
\iqwhy{This is why fast Fourier transforms solve circular convolutions: the Fourier basis diagonalises every circulant matrix at once.}
\end{solution}

\begin{solution}{iq:iv1:linear-algebra:8}
For $n = 3$ the eigenvalues are $2 - \sqrt{2}$, $2$ and $2 + \sqrt{2}$. In general they are $2 - 2\cos\frac{k\pi}{n + 1}$ for $k = 1, \dots, n$. The condition number grows like $4(n + 1)^2/\pi^2$, about 4100 for $n = 100$.

Try $x_j = \sin\frac{jk\pi}{n + 1}$, a discrete sine wave that is zero at $j = 0$ and $j = n + 1$\index{second-difference matrix}. Then
\[
-x_{j - 1} + 2x_j - x_{j + 1} = \Bigl(2 - 2\cos\frac{k\pi}{n + 1}\Bigr)x_j,
\]
by the sum formula for sines. It satisfies the boundary rows automatically.

The largest eigenvalue is near 4, and the smallest is about $\pi^2/(n + 1)^2$, so the ratio goes as stated. A finer grid gives a less well-conditioned matrix. Direct solves are still fine for a tridiagonal system, but rounding errors grow with the grid, and an iterative solver slows down without a preconditioner.

\iqwrong{Thinking that a finer grid must always give better answers. The discretisation error falls, but the rounding error rises, and there is a best grid size.}
\end{solution}

\begin{solution}{iq:iv1:linear-algebra:9}
With unit volatilities the components are $(1, 1)/\sqrt{2}$, with a 75\% share, and $(1, -1)/\sqrt{2}$, with 25\%. If the first asset's volatility is doubled, the first component explains about 86\%. PCA depends on units, so use the correlation matrix when units are not comparable.

The first matrix has eigenvalues $1 + \rho = 1.5$ and $1 - \rho = 0.5$\index{principal component analysis}, summing to 2.

With volatilities 2 and 1, the covariance matrix has rows $(4, 1)$ and $(1, 1)$. Its eigenvalues are $\frac{5 \pm \sqrt{13}}{2}$, so $4.30$ and $0.70$. The share is $4.30/5 = 0.86$. The top component now points mostly along the first asset, at about 17 degrees from it, with $\tan 2\theta = 2/3$. The correlation has not changed. Only the units have.

So PCA on returns with very different volatilities mostly finds the most volatile assets. Decide whether you want that. Scale to unit variance if not, or weight by a risk budget.

\iqwrong{Treating the components as properties of the assets. They are also a property of the scale you chose, and a change of currency or a leverage factor changes them.}
\end{solution}

\begin{solution}{iq:iv1:linear-algebra:10}
The singular values are $3\sqrt{5} \approx 6.71$ and $\sqrt{5} \approx 2.24$, so the condition number is 3. The eigenvalues are 3 and 5, which would suggest a different ratio.

Singular values are the square roots of the eigenvalues of $A^\top A$\index{singular value decomposition}. Here $A^\top A$ has rows $(25, 20)$ and $(20, 25)$, whose eigenvalues are 45 and 5.

The matrix is triangular, so its eigenvalues are its diagonal: 3 and 5. For a non-symmetric matrix, eigenvalues do not measure size or conditioning, and only singular values do.

For a centred data matrix $X = U\Sigma V^\top$ with $n$ rows, the covariance is $X^\top X/(n - 1) = V\Sigma^2V^\top/(n - 1)$. So the columns of $V$ are the principal components, and their variances are $\sigma_i^2/(n - 1)$. A test on random data confirms this. It is better to use the SVD on $X$ than an eigensolver on the covariance, because forming the covariance squares the condition number.

\iqwrong{Using the ratio $5/3 \approx 1.67$ as the condition number. That works only for a symmetric matrix.}
\end{solution}

\begin{solution}{iq:iv1:linear-algebra:11}
The solutions are $(2, 0)$ and $(1, 1)$. So changing $b$ by $0.0001$ moves $x$ by 1. The condition number is about $4 \times 10^4$, and you lose about $4.6$ digits. Forming $A^\top A$ squares the condition number, which loses about $9$ digits.

The matrix is nearly singular\index{condition number}, as its rows are nearly equal. Its eigenvalues are about 2 and $5 \times 10^{-5}$, whose ratio is the condition number.

The rule of thumb is that relative error in the answer can be $\kappa$ times the relative error in the input. If you work with 16 digits, you can lose about $\log_{10}\kappa$ of them. With the normal equations, $X^\top X$ has a condition number of $\kappa^2$, and that is why good regression code uses QR or the SVD on $X$.

\iqwrong{Checking the residual, and being satisfied that it is small. Both solutions have tiny residuals against the other $b$. A small residual does not mean a small error when $A$ is badly conditioned.}
\iqwhy{Two near-identical assets in a risk model create exactly this problem. A small change in estimated means or covariances then gives wildly different portfolio weights.}
\end{solution}

\begin{solution}{iq:iv1:linear-algebra:12}
The rank is at most 249, after subtracting the means. So there are at least 251 independent portfolios with exactly zero sample variance, and a minimum-variance optimiser sends its weights to them.

$S = \frac{1}{T - 1}\sum_t (r_t - \bar{r})(r_t - \bar{r})^\top$ is a sum of $T$ rank-one matrices\index{sample covariance matrix}, and the centred vectors add up to zero, so its rank is at most $T - 1$. A simulation with 250 and 500 confirms 249.

A portfolio with zero sample variance is not riskless. Its weights are fitted to the sample's quirks, and its true risk is ordinary. The optimiser cannot be solved uniquely and, with constraints, it still heads straight for these directions.

Two repairs: impose structure with a factor model, $\Sigma = D + BB^\top$, which is full rank with few parameters. Or shrink: $\alpha F + (1 - \alpha)S$, with a simple target $F$. Weight limits also help, because they act like shrinkage.

\iqwrong{Inverting with a pseudo-inverse and calling it fixed. That gives you a number, but the problem is that you have too little data, and no matrix algebra adds any.}
\end{solution}

\begin{solution}{iq:iv1:linear-algebra:13}
The new eigenvalues are $4.11$, $1.31$ and $0.68$. The condition number falls from 50 to about $6.0$. The eigenvectors and the trace stay the same.

With $\mu = 6.1/3 \approx 2.03$\index{shrinkage estimators}, each eigenvalue becomes $0.7\lambda + 0.3\mu$, because $I$ has every vector as an eigenvector. Sample eigenvalues are too spread out: the top is too high, and the bottom too low. So this moves them together while keeping the total variance.

An optimiser's weights use $\Sigma^{-1}$, so they are dominated by the smallest eigenvalues, which are the least reliable. Raising them from $0.1$ to $0.68$ reduces the weight given to those directions about sevenfold.

The $0.3$ can be chosen by the Ledoit--Wolf formula, which estimates the intensity that minimises the expected squared distance to the true matrix, using how noisy $S$ is. Or use cross-validation on out-of-sample variance.

\iqwrong{Thinking of shrinkage as losing information because it changes the numbers. The sample values have errors that always spread them apart, and moving them together corrects that.}
\end{solution}

\begin{solution}{iq:iv1:linear-algebra:14}
The weights are proportional to $\sigma_2^2 - \rho\sigma_1\sigma_2$ and $\sigma_1^2 - \rho\sigma_1\sigma_2$. The second is zero at $\rho = \sigma_1/\sigma_2 = 0.5$. At $\rho = 0.8$ the weights are $(4/3, -1/3)$, and the volatility is about $8.94\%$.

$w \propto \Sigma^{-1}\mathbf{1}$\index{minimum-variance portfolio}, which for two assets gives the formula above. A grid search on the weight at $\rho = 0.8$ finds the same.

At $\rho = 0.8$: variance is $\frac{16}{9}(0.01) + \frac19(0.04) - 2 \cdot \frac43 \cdot \frac13 \cdot 0.016 = 0.008$.

Why the variance is lower: the short position in the second asset hedges the first. The marginal effect of adding the second asset at zero weight is its covariance with the first, $0.016$, which is above the first's variance, $0.01$. So being short it lowers the risk. In other words, the beta of asset 1 on asset 2 is $0.4$, and the best portfolio holds a partial hedge against it.

Above $\rho = 0.5$ the riskier asset is a hedge, and it is shorted. Below that, you hold it long.

\iqwrong{Saying that a riskier asset always gets a smaller positive weight. For a high enough correlation it gets a negative weight, because it is then useful as a hedge.}
\end{solution}

\begin{solution}{iq:iv1:linear-algebra:15}
The matrix has rows $(0.9, 0.1)$ and $(0.2, 0.8)$, with eigenvalues 1 and $0.7$. The stationary distribution is $(2/3, 1/3)$. The effect of today's state shrinks by a factor of $0.7$ per day, which is a half-life of about $1.9$ days.

Rows sum to 1, so $\mathbf{1}$ is a right eigenvector with eigenvalue 1\index{Markov chain eigenvalues}. The other eigenvalue comes from the trace: $1.7 - 1 = 0.7$. The stationary $\pi$ solves $\pi P = \pi$, which gives $0.1\pi_1 = 0.2\pi_2$, so $\pi_1 = 2\pi_2$.

The difference between the chance of calm tomorrow starting from calm and starting from stress is $0.9 - 0.2 = 0.7$, which is the eigenvalue. After $k$ days that is $0.7^k$, so half-way takes $\ln 0.5/\ln 0.7 \approx 1.94$ days. A matrix power at 200 steps gives the stationary row from either start.

\iqwrong{Thinking the state's duration tells you about mixing: the expected length of a stress spell is 5 days and of a calm spell is 10, and these are longer than the half-life. Mixing combines both exits, and the eigenvalue already includes them.}
\end{solution}

\begin{solution}{iq:iv1:linear-algebra:16}
The matrix has rows $(1, 1)$ and $(1, 0)$, with eigenvalues $\varphi = (1 + \sqrt{5})/2$ and $\psi = (1 - \sqrt{5})/2$. So $F_n = (\varphi^n - \psi^n)/\sqrt{5}$, and about $2\log_2 n$ matrix products are needed. With doubles, the rounded formula first goes wrong at $n = 71$.

$(F_{n + 1}, F_n)^\top = M(F_n, F_{n - 1})^\top$, so $M^n$ has $F_n$ in its off-diagonal\index{matrix powers}. Expanding the start vector in eigenvectors gives the formula. Since $|\psi| < 1$, $F_n$ is the nearest integer to $\varphi^n/\sqrt{5}$.

Repeated squaring computes $M^n$ in about $\log_2 n$ squarings, and up to as many extra products.

Doubles have about 53 bits. $\varphi^n/\sqrt{5}$ passes $2^{53}$ near $n = 78$, but the rounding errors in $\varphi$ accumulate before then. A test finds the first wrong value at $n = 71$. So use integer matrix powers for an exact value. A test computes $F_{30} = 832\,040$ that way.

\iqwrong{Saying that the closed form is always better because it has no loop. It is only as exact as its floating-point inputs, and the integer method has no such limit.}
\end{solution}

\begin{solution}{iq:iv1:linear-algebra:17}
The rank is 2, and the nonzero eigenvalues are $5050 \pm \sqrt{100 \cdot 338\,350} \approx 10\,866.8$ and $-766.8$.

Write $A = u\mathbf{1}^\top + \mathbf{1}u^\top$\index{low-rank structure}, with $u = (1, \dots, 100)$. Each term has rank 1, and the two together have rank 2. Any eigenvector with a nonzero eigenvalue lies in the span of $u$ and $\mathbf{1}$. Write $x = au + b\mathbf{1}$. Then $Ax = u(\mathbf{1}^\top x) + \mathbf{1}(u^\top x)$, which gives the 2 by 2 matrix
\[
\begin{pmatrix} \mathbf{1}^\top u & \mathbf{1}^\top\mathbf{1} \\ u^\top u & u^\top\mathbf{1} \end{pmatrix} = \begin{pmatrix} 5050 & 100 \\ 338\,350 & 5050 \end{pmatrix}.
\]
Its eigenvalues are $5050 \pm \sqrt{100 \cdot 338\,350}$. As a check, the sum is $10\,100$, which is the trace of $A$: $\sum 2i = 10\,100$. A numerical solver gives the same.

\iqwrong{Expecting all eigenvalues to be positive, because all entries are positive. The matrix is not positive semi-definite. A positive entry does not imply that.}
\end{solution}

\begin{solution}{iq:iv1:linear-algebra:18}
The sample eigenvalues lie between about $0.086$ and $2.914$. In the sample, the portfolio shows a variance of $(1 - q)/N$, with $q = N/T = 0.5$. Out of sample, its true variance is $1/((1 - q)N)$. So the optimum variance is $1/N$, the sample reports half of that, and it delivers twice that, for a ratio of 4 between the delivered risk and the risk seen.

The law says that the eigenvalues of pure noise spread out to $(1 \pm \sqrt{q})^2$\index{Marchenko--Pastur law}, so $(1 \pm 0.707)^2$. A simulation with 500 assets and 1000 days gets close to both edges.

The optimiser takes $w \propto S^{-1}\mathbf{1}$. So it picks up the low sample eigenvalues, which are spuriously low, and loads onto those directions. Twenty simulations, one for each seed, give in-sample variance $0.5/N$ and out-of-sample variance $2/N$, against the true optimum $1/N$, which is the equal-weight portfolio here.

To make this useful, look at the limits. As $q \to 0$, you have lots of data, and all three agree. As $q \to 1$, the out-of-sample risk blows up. Below the Marchenko--Pastur upper edge, real data's eigenvalues cannot be told apart from noise, so clean them or use a factor model.

\iqwrong{Thinking that the risk in the sample is a fair estimate of the risk you will face. An optimised portfolio is the most flattering one for that sample, so it is too low by construction.}
\end{solution}

\begin{solution}{iq:iv1:linear-algebra:19}
$\Sigma^{-1} = D^{-1} - \dfrac{D^{-1}\beta\beta^\top D^{-1}}{1 + \beta^\top D^{-1}\beta}$. This takes $O(N)$ operations for a matrix-vector product, instead of $O(N^3)$ for a direct inverse. The minimum-variance weights are proportional to $d_i^{-1}(1 - c\beta_i)$, so high-beta stocks get low or negative weights.

Proof\index{Sherman--Morrison formula}: multiply out. Let $x = D^{-1}\beta$ and $s = \beta^\top x$. Then
\[
(D + \beta\beta^\top)\Bigl(D^{-1} - \frac{xx^\top}{1 + s}\Bigr) = I + \beta x^\top - \frac{\beta x^\top + \beta s x^\top}{1 + s} = I + \beta x^\top - \beta x^\top = I.
\]
Numerical tests with 1000 assets confirm this, and the fast weights match a direct solve.

Applying it to $\mathbf{1}$ gives $w \propto D^{-1}\mathbf{1} - x\frac{\mathbf{1}^\top x}{1 + s}$, whose $i$-th entry is $\frac{1}{d_i}\bigl(1 - \beta_i\frac{\mathbf{1}^\top x}{1 + s}\bigr)$, which is the form above. So the weights fall in a straight line with beta, when residual risks are equal. That is the known low-beta tilt of minimum-variance portfolios: the portfolio lowers its market exposure by buying low-beta and shorting high-beta names.

The cost matters: with $k$ factors, Woodbury's version needs the inverse of one $k \times k$ matrix, at a cost of $O(Nk^2)$. So factor models are as much a computational tool as a statistical one.

\iqwrong{Forming and inverting the whole matrix each time. That costs a billion operations at this size, and it loses the structure that makes the result stable.}
\end{solution}

\begin{solution}{iq:iv1:linear-algebra:20}
The partial correlation is $1/3$. In general it is $\rho/(1 + \rho)$.

From residuals: regressing on $Z$ gives $X - 0.5Z$ and $Y - 0.5Z$\index{partial correlation}. Each has variance $1 - 0.25 = 0.75$. Their covariance is $0.5 - 0.25 - 0.25 + 0.25 = 0.25$. So the correlation is $0.25/0.75 = 1/3$. A simulation agrees. In general, $\frac{\rho_{XY} - \rho_{XZ}\rho_{YZ}}{\sqrt{(1 - \rho_{XZ}^2)(1 - \rho_{YZ}^2)}}$.

From the inverse: let $P = C^{-1}$. Then the partial correlation is $-P_{12}/\sqrt{P_{11}P_{22}}$\index{precision matrix}. For this equicorrelated matrix, the inverse has diagonal $\frac{1 + \rho}{(1 - \rho)(1 + 2\rho)}$ and off-diagonal $\frac{-\rho}{(1 - \rho)(1 + 2\rho)}$, which gives $\rho/(1 + \rho)$.

So the off-diagonal entries of the inverse measure direct links, once all the other variables have been accounted for. A zero entry means that two variables are conditionally independent, in the Gaussian case.

Limits: $\rho = 0$ gives 0, and $\rho \to 1$ gives $\tfrac12$. It is not 1, because when all are close to identical, the residuals after removing $Z$ are still correlated by the common structure.

\iqwrong{Using the correlation $0.5$ as the direct link. Part of it comes through $Z$. Partial correlation removes that part, and it may fall a lot, or even change sign.}
\end{solution}

\begin{solution}{iq:iv1:linear-algebra:21}
Multiply a start vector by the matrix, then normalise, and repeat. It takes about 20 steps for a $10^{-6}$ error in the vector, and 10 for the eigenvalue. If the second eigenvalue is $9.9$, it takes about 1375 steps. It fails if the start vector has no component along the top vector.

Write the start as $\sum c_iv_i$\index{power iteration}. After $k$ steps, it is $\lambda_1^k\bigl(c_1v_1 + c_2(\lambda_2/\lambda_1)^kv_2 + \dots\bigr)$. So the error falls like $(\lambda_2/\lambda_1)^k$, and $0.5^{20} \approx 10^{-6}$.

The Rayleigh quotient $x^\top Ax$ has an error that depends on the square of the vector's error, for a symmetric matrix, so it needs half as many steps. A run on a 1000 by 1000 matrix confirms both orders of size.

For $9.9/10$, $k = \ln 10^{-6}/\ln 0.99 \approx 1375$. So the gap matters, not the size.

Failure: exact orthogonality to $v_1$, which is unlikely from a random start, but rounding would usually fix it by itself anyway. A test that keeps removing the top direction shows it stuck on the other direction. If $\lambda_1 = \lambda_2$ the top eigenvector is not unique: the iteration converges to some vector in that plane, and which one depends on the start. Lanczos is far faster for a small gap, using the same matrix-vector product, and needs no forming of the matrix. That matters when the covariance is $XX^\top$ or a factor model.

\iqwrong{Giving an answer that depends on the matrix's size. The cost per step does, through the matrix-vector product. The number of steps depends only on the ratio of the eigenvalues.}
\end{solution}

\begin{solution}{iq:iv1:linear-algebra:22}
If $ABv = \lambda v$ with $\lambda \ne 0$, then $w = Bv$ is non-zero, and $BAw = BABv = \lambda Bv = \lambda w$. So nonzero eigenvalues are shared, with the same multiplicities. For the returns, form the 250 by 250 matrix $XX^\top$ and map each eigenvector back with $v = X^\top u/\sqrt{\lambda}$.

With $A = X$ and $B = X^\top$\index{dual PCA}, $XX^\top$ and $X^\top X$ have the same nonzero eigenvalues. If $XX^\top u = \lambda u$ and $u$ has unit length, let $v = X^\top u/\sqrt{\lambda}$. Then $|v|^2 = u^\top XX^\top u/\lambda = 1$, and $X^\top Xv = X^\top XX^\top u/\sqrt{\lambda} = \lambda X^\top u/\sqrt{\lambda} = \lambda v$. A run with $T = 250$ and $N = 5000$ confirms this. The eigenproblem takes about $(250/5000)^3 \approx 10^{-4}$ of the work.

For the determinant, the identity $\det(I + AB) = \det(I + BA)$ follows from the shared nonzero eigenvalues: each $1 + \lambda$ appears in both, and the rest are 1. With $A = u$ and $B = v^\top$, $BA$ is the 1 by 1 number $v^\top u$. So $\det(I + uv^\top) = 1 + v^\top u$. Alternatively, $uv^\top$ has one nonzero eigenvalue $v^\top u$, and the others are 0.

\iqwrong{Forming the $5000 \times 5000$ matrix first because that is what the formula says. Its rank is at most 249, so most of that work goes on zeros.}
\end{solution}

\begin{solution}{iq:iv1:linear-algebra:23}
The eigenvalues are $1 + 0.9\sqrt{2} \approx 2.273$, $1$ and $1 - 0.9\sqrt{2} \approx -0.273$. After clipping and rescaling, $\rho_{12} = \rho_{23} \approx 0.729$ and $\rho_{13} \approx 0.064$. This is close, but it is not the closest valid correlation matrix: that one has about $0.74$ and $0.105$.

Set the negative eigenvalue to 0\index{nearest correlation matrix}. Let $v = (\tfrac12, -\tfrac{1}{\sqrt{2}}, \tfrac12)$ be its eigenvector. Then
\[
A' = A + 0.273\,vv^\top,
\]
which has diagonal $(1.068, 1.136, 1.068)$. It has off-diagonals $0.900 - 0.273 \cdot 0.354 \approx 0.803$, and $0.273 \cdot \tfrac14 \approx 0.068$. Dividing by the square roots of the diagonals gives the new correlations.

In what sense is it close: before the rescaling, $A'$ is the nearest positive semi-definite matrix in the Frobenius norm, at a distance of $0.273$, which is exactly the clipped eigenvalue. Rescaling then moves it away, because it restores the unit diagonal without thinking about distance.

The nearest correlation matrix alternates projections onto the positive semi-definite cone and onto unit diagonals, with Dykstra's correction. A test with this method gives $0.74$ and $0.105$, a smaller distance from the original matrix. The two differ little in many uses, and clipping is simple and fast.

\iqwrong{Assuming that rescaling after clipping is the nearest matrix, or that you can fix the problem by adjusting the worst single number. Any repair changes several entries at once.}
\end{solution}

\begin{solution}{iq:iv1:linear-algebra:24}
$E[\det(A)^2] = n!$. For $n = 3$ with $\pm 1$ entries, there are $2^9 = 512$ matrices, and the average of $\det^2$ is 6.

Expand the determinant by permutations\index{random determinants}:
\[
\det(A) = \sum_\sigma \operatorname{sgn}(\sigma)\prod_i a_{i\sigma(i)}.
\]
Squaring gives a double sum over $\sigma$ and $\tau$. If $\sigma \ne \tau$, there is a row $i$ where $\sigma(i) \ne \tau(i)$, so $a_{i\sigma(i)}$ appears alone in that product (each entry only once in each of the two products), and it has mean zero. If $\sigma = \tau$, the product is $\prod a_{i\sigma(i)}^2$, which has mean 1. So only the $n!$ diagonal terms survive.

The enumeration of all 512 matrices gives exactly 6. Simulations with normal entries for $n = 2$, 4 and 5 agree. So the formula needs only independence, mean 0 and variance 1.

So the root-mean-square size is $\sqrt{n!}$, which sets the growth rate. A typical value is smaller: with normal entries, the geometric mean falls short by a factor of about $1.7\sqrt{n}$, little beside the factorial. For $\pm 1$ matrices, the largest possible is at most $n^{n/2}$, by Hadamard's inequality\index{Hadamard's inequality}, because the determinant is the volume of a box with sides of length $\sqrt{n}$. For $n = 3$ the bound is $5.2$, and the enumeration finds that the actual maximum is 4. The ratio between $\sqrt{n!}$ and $n^{n/2}$ falls like $e^{-n/2}$, by Stirling. So random vectors are far from orthogonal in that volume sense, even though each pair is nearly orthogonal.

\iqwrong{Saying that the mean of the determinant is 0, and so the typical size is small. Yes, the mean is 0, since every term of the expansion is a product of independent factors with mean 0. The second moment gives its size, and it is huge.}
\end{solution}

\begin{solution}{iq:iv1:linear-algebra:25}
Write $A = \sum_k \lambda_ku_ku_k^\top$ and $B = \sum_l \mu_lv_lv_l^\top$, with all coefficients non-negative. Then $A \circ B = \sum_{k, l}\lambda_k\mu_l(u_k \circ v_l)(u_k \circ v_l)^\top$, which is a sum of positive semi-definite matrices\index{Schur product theorem}.

To check the middle step: entry $(i, j)$ of $A \circ B$ is $\bigl(\sum_k\lambda_ku_{ki}u_{kj}\bigr)\bigl(\sum_l\mu_lv_{li}v_{lj}\bigr)$, which is the same as $\sum_{k, l}\lambda_k\mu_l(u_{ki}v_{li})(u_{kj}v_{lj})$. Random tests with low-rank matrices never find a negative eigenvalue, beyond rounding.

For the taper, $e^{-\lambda|i - j|}$ is the correlation matrix of an AR(1) process, and so it is positive definite. So the product with a correlation matrix is positive semi-definite. The diagonal stays 1, because both diagonals are 1. It is a valid correlation matrix\index{covariance tapering}. A test with $n = 8$ confirms this.

For the bound, let $m = \lambda_{\min}(A)$. Then $A - mI$ is positive semi-definite, so $(A - mI) \circ B$ is too. This expands to $A \circ B - m\operatorname{diag}(B)$. So
\[
x^\top(A \circ B)x \ge m\sum_i B_{ii}x_i^2 \ge m\min_iB_{ii}|x|^2.
\]
Random tests confirm this inequality too.

This is why a smooth taper is safe, and a hard cut-off is not. Cutting all entries beyond a band to zero is an element-wise product with a matrix that is not positive semi-definite, and a test shows that it can yield a negative eigenvalue: the one-factor matrix with $\rho = 0.9$ and a band of width 1 gives one.

\iqwrong{Trying to use the eigenvalues of the two matrices directly, as with a matrix product. The eigenvalues of $A \circ B$ are not products of eigenvalues. The decomposition into rank-one terms is what makes it work.}
\end{solution}
